%! TeX program = lualatex
\documentclass[../main.tex]{subfiles}
\begin{document} \section{The chain rule\texorpdfstring{ \((f \circ g)' = f'(g(x))g'(x)\)}{}}

This section is about building procedural competency for the chain rule.

\begin{mdframed}[style=simple-compact]
  \textbf{Differentiation Rules} \hfill {\footnotesize (part \(3\) of \(3\))}

  The derivative of a composite function \(f(g(x))\), if it is well-defined, can be calculated by the formula \updatenote{Added the missing "it".}
  \begin{align*}
    \frac{d}{dx} f\big( g(x) \big)
    &= f'(g(x)) g'(x)
    &&\text{(the chain rule)}\\[1em]
    \frac{dy}{dx} &= \frac{dy}{du} \frac{du}{dx} \text{ if we write } y = f(u), u = g(x).
  \end{align*}
\end{mdframed}

\begin{example}
  Differentiate \(\tan^{3}(x)\).

  \blanklines{10}
\end{example}

\begin{exercise}
  Differentiate \(e^{2x + 1}\).

  \blanklines{10}
\end{exercise}

\begin{exercise}
  Differentiate \(\left( x - \frac{4}{x} \right)^{8}\).

  \blanklines{10}
\end{exercise}

\clearpage

\begin{exercise}
  Find all \(x\) at which the curve \(y = \left( x - \frac{4}{x} \right)^{8}\) has a horizontal tangent line.

  \blanklines{15}
\end{exercise}

The chain rule can help us differentiate composition functions where part or all of the functions are unknown.
\begin{example}
  Differentiate \(\sec(g(x))\) where \(g\) is an unknown but differentiable function. 

  \blanklines{5}
\end{example}

\begin{exercise}
  Differentiate \((4 g(x) + 3)^{3}\) where \(g(x)\) is an unknown but differentiable function.

  \blanklines{10}
\end{exercise}

\begin{exercise}
  Find the derivative of \(\sqrt{2^{x} + g(x)}\) at \(x = 1\) given that \(g(1) = 1\) and \(g'(1) = 0\).

  \blanklines{15}
\end{exercise}

\begin{exercise}
  Find all \(x\) at which \(\cos^{2}(x)\) has a horizontal tangent line.

  \blanklines{20}
\end{exercise}

\begin{exercise}
  Find all \(x\) at which \(\cos^{2}(x)\) has a tangent line with slope \(1\).

  \blanklines{20}
\end{exercise}
\clearpage
\end{document}
